\documentclass[../../../include/open-logic-section]{subfiles}

\begin{document}

\olfileid[hi]{sth}{story}{predicative}	

\olsection{विधेयात्मक और अविधेयात्मक}

रसेल समुच्चय $R$ को $\Setabs{x}{x \notin x}$ द्वारा परिभाषित किया गया
था। अधिक विस्तार से कहें तो $R$ वह समुच्चय होता जिसमें वे सभी और केवल वे
समुच्चय होते जो स्वयं के असदस्य नहीं हैं। अतः $R$ को परिभाषित करते समय हम
उस प्रांत पर परिमाणीकरण करते हैं जिसमें $R$ भी होता (यदि उसका अस्तित्व
होता)।

यह एक \emph{अविधेयात्मक} परिभाषा है। अधिक सामान्यतः हम कह सकते हैं कि कोई
परिभाषा अविधेयात्मक तभी और केवल तभी है जब वह ऐसे प्रांत पर परिमाणीकरण करती
है जिसमें परिभाषित की जा रही वस्तु सम्मिलित हो।
	
विरोधाभासों के सामने आने के बाद व्हाइटहेड, रसेल, प्वाँकारे और वाइल ने ऐसी
अविधेयात्मक परिभाषाओं को ``दूषित रूप से चक्रीय'' कहकर अस्वीकार किया:
\begin{quote}
	जिन विरोधाभासों से बचना है उनका विश्लेषण दिखाता है कि वे सभी एक प्रकार
	के दूषित चक्र से उत्पन्न होते हैं। ये दूषित चक्र इस धारणा से पैदा होते
	हैं कि वस्तुओं के किसी संग्रह में ऐसे सदस्य हो सकते हैं जिन्हें केवल
	संपूर्ण संग्रह के माध्यम से ही परिभाषित किया जा सकता है[\ldots.
	\textparagraph]

	अवैध समग्रताओं से बचाने वाला सिद्धांत इस प्रकार कहा जा सकता है: `जिसमें
	किसी संग्रह के \emph{सभी} सदस्य अंतर्भूत हों, वह स्वयं उस संग्रह का एक
	सदस्य नहीं होना चाहिए'; अथवा इसके विपरीत: `यदि किसी संग्रह की समग्रता
	मान लेने पर उसमें ऐसे सदस्य होते जिन्हें केवल उस समग्रता के पदों में
	परिभाषित किया जा सकता, तो उस संग्रह की कोई समग्रता नहीं है।' हम इसे
	`दूषित-चक्र सिद्धांत' कहेंगे, क्योंकि यह हमें अवैध समग्रताओं की धारणा में
	निहित दूषित चक्रों से बचने देता है।
	\citep[p.~37]{WhiteheadRussell1910}
\end{quote}
यदि हम उनके साथ चलते हुए \emph{दूषित-चक्र सिद्धांत} स्वीकार करें, तो
\olref[sth][story][rus]{sec} की विनाशकारी सहज समुच्चय-निर्माण योजना को कुछ
इस प्रकार बदलने का प्रयास कर सकते हैं:

\begin{defish}
\emph{विधेयात्मक समुच्चय-निर्माण।} केवल समुच्चयों पर परिमाणीकरण करने वाले
प्रत्येक सूत्र $\phi$ के लिए: समुच्चय$^\prime$
$\Setabs{x}{\phi(x)}$ का अस्तित्व है।
\end{defish}

जब तक समुच्चय$^{\prime}$ समुच्चय नहीं हैं, कोई विरोधाभास उत्पन्न नहीं होगा।

दुर्भाग्य से विधेयात्मक समुच्चय-निर्माण बहुत \emph{व्यापक} नहीं है। आखिर
यह हमें नई सत्ताओं, समुच्चयों$^\prime$, से परिचित कराता है। इसलिए हमें उन
सूत्रों पर विचार करना होगा जो समुच्चयों$^\prime$ पर परिमाणीकरण करते हैं। यदि
वे सदा कोई समुच्चय$^\prime$ देते हैं, तो समस्त स्वयं के असदस्य
समुच्चयों$^\prime$ के समुच्चय$^\prime$ पर विचार करते ही रसेल का विरोधाभास
फिर उत्पन्न होगा। अतः इसी विचार को आगे बढ़ाते हुए हमें कहना होगा कि
समुच्चयों$^\prime$ पर परिमाणीकरण करने वाला सूत्र कोई संगत
समुच्चय$^{\prime\prime}$ देता है। फिर हमें समुच्चय$^{\prime\prime\prime}$,
समुच्चय$^{\prime\prime\prime\prime}$ इत्यादि चाहिए होंगे। प्राइम चिह्नों की
बाढ़ रोकने के लिए इन्हें समुच्चय$_0$, समुच्चय$_1$, समुच्चय$_2$,
समुच्चय$_3$, समुच्चय$_4$,\ldots{} मानना सरल होगा। इससे हमें (सरल) प्रकार
सिद्धांत में प्रवेश का एक मार्ग मिलेगा।

ऐसे सिद्धांत के विरुद्ध कुछ स्पष्ट आपत्तियाँ हैं (यद्यपि यह स्पष्ट नहीं कि
वे \emph{निर्णायक} आपत्तियाँ हैं)। संक्षेप में: परिणामी सिद्धांत का उपयोग
दुष्कर है; वह भिन्न-भिन्न प्रकार की वस्तुएँ मानने में अपव्ययी है; और अंततः
यह भी स्पष्ट नहीं कि अविधेयात्मक परिभाषाएँ सचमुच \emph{इतनी बुरी} हैं।
	
अंतिम बात स्पष्ट करने के लिए \citeauthor{Ramsey1925} की यह टिप्पणी देखें:
\begin{quote}
	हम किसी व्यक्ति को किसी समूह का सबसे लंबा व्यक्ति कह सकते हैं और इस तरह
	उसे ऐसी समग्रता के माध्यम से पहचान सकते हैं जिसका वह स्वयं सदस्य है,
	फिर भी इसमें कोई दूषित चक्र नहीं होता। \citep{Ramsey1925}
\end{quote}
रैम्ज़ी की बात यह है कि ``समूह का सबसे लंबा व्यक्ति'' एक अविधेयात्मक
परिभाषा \emph{है}; किंतु वह स्पष्टतः पूरी तरह निर्दोष है।

उत्तर दिया जा सकता है कि इस मामले में सबसे लंबे व्यक्ति को
\emph{विधेयात्मक} साधनों से चुना जा सकता है। उदाहरणतः शायद हम बस उस व्यक्ति
की ओर संकेत कर दें। तब अविधेयात्मक परिभाषाओं के विरुद्ध आपत्ति स्पष्टतः उन
सत्ताओं तक सीमित करनी होगी जिन्हें \emph{केवल} अविधेयात्मक ढंग से चुना जा
सकता है। पर तब भी हमें यह और समझाना होगा कि ऐसी ``अनिवार्य अविधेयात्मकता''
इतनी बुरी क्यों होगी।\footnote{और अधिक के लिए \citet{Linnebo2010} देखें।}

यह स्वीकार करना होगा कि यदि हम चाहते हैं कि हमारी परिभाषाएँ किसी वस्तु को
\emph{बनाने} की विधि जैसी कोई चीज़ दें, तो अविधेयात्मक परिभाषाएँ अत्यंत
समस्याजनक हैं। किसी अविधेयात्मक परिभाषा के साथ हम सचमुच दूषित चक्र में फँस
जाएँगे: अविधेयात्मक रूप से निर्दिष्ट वस्तु बनाने के लिए पहले सभी वस्तुएँ
(स्वयं उस अविधेयात्मक रूप से निर्दिष्ट वस्तु सहित) बनानी होंगी, क्योंकि उसका
निर्देश सभी वस्तुओं के पदों में होता है; अतः उसे स्वयं बनाने से पहले ही
बनाना पड़ेगा। किंतु फिर, ``अनिवार्य रूप से अविधेयात्मक'' ढंग से निर्दिष्ट
समुच्चयों के विरुद्ध यह गंभीर आपत्ति तभी है जब हम समुच्चयों को ऐसी चीज़ें
मानें जिन्हें हम \emph{बनाते} हैं। और हम (संभवतः) ऐसा नहीं मानते।

इसलिए---भला हो या बुरा---जो दृष्टिकोण प्रचलित हुआ, वह
(अ)\-वि\-धे\-यात्\-मक\-ता के विषय में कठोर रुख नहीं अपनाता। इसके बजाय उसमें
वह दृष्टिकोण है जिसे अब संचयी-पुनरावर्ती दृष्टिकोण माना जाता है। अंततः यह
हमें अपने समुच्चयों को ``चरणों'' में स्तरित करने देगा---कुछ-कुछ वैसे ही जैसे
विधेयात्मक दृष्टिकोण सत्ताओं को समुच्चय$_0$, समुच्चय$_1$, समुच्चय$_2$,
\ldots{} में स्तरित करता है---पर हम उनके बीच प्रकार का कोई भेद नहीं मानेंगे।

\end{document}
